\documentclass[11pt]{article}
\usepackage[a5paper,margin=18mm]{geometry}
\usepackage{amsmath}
\usepackage{mathtools}
\usepackage{unicode-math}
\setmainfont{TeX Gyre Pagella}
\setsansfont{TeX Gyre Heros}
\setmonofont{TeX Gyre Cursor}
\setmathfont{texgyrepagella-math.otf}
\setmathfont{latinmodern-math.otf}[range={cal,bfcal},StylisticSet=1]
\DeclarePairedDelimiter{\abs}{\lvert}{\rvert}
\DeclarePairedDelimiter{\norm}{\lVert}{\rVert}
\DeclareMathOperator{\Tr}{Tr}

\begin{document}
\section*{Spectral identities}

Let $A\in\mathbb{R}^{2\times 2}$ be symmetric with eigenvalues $\lambda_1,\lambda_2$.
Then, with $\mu \coloneqq \tfrac12\Tr A$ and $\delta \coloneqq \tfrac12\abs{\lambda_1-\lambda_2}$,
\begin{align}
  \lambda_{1,2} &= \mu \pm \delta, \label{eq:eig}\\
  \Tr A^{2} &= \lambda_1^{2}+\lambda_2^{2} = 2\mu^{2}+2\delta^{2}. \label{eq:tr}
\end{align}
Equation~\eqref{eq:eig} fixes the spectrum, and \eqref{eq:tr} its second moment, hence
\begin{equation}
  \norm{A}_{F}^{2} = \sum_{\mathclap{i,j=1}}^{2} a_{ij}^{2}
  = \underbrace{\lambda_1^{2}+\lambda_2^{2}}_{\text{trace of square}}.
  \label{eq:frob}
\end{equation}

\subsection*{A piecewise definition}
\[
  \operatorname{sgn}(x)=
  \begin{dcases}
    -1 & \text{if } x<0,\\
    \phantom{-}0 & \text{if } x=0,\\
    +1 & \text{if } x>0.
  \end{dcases}
\]

\subsection*{Matrices and arrows}
\[
  \mathbf{R}(\theta)=\begin{pmatrix}\cos\theta & -\sin\theta\\ \sin\theta & \phantom{-}\cos\theta\end{pmatrix}
  \xrightarrow[\text{rotate by }\pi]{\theta\to\pi}
  \begin{bmatrix}-1&0\\0&-1\end{bmatrix}
\]
Blackboard $\mathbb{N}\subset\mathbb{Z}\subset\mathbb{Q}\subset\mathbb{R}\subset\mathbb{C}$,
calligraphic $\mathcal{L}\mathcal{F}$, bold italic $\symbfit{x}$,
and the Greek set $\alpha\beta\gamma\Gamma\Delta\Omega\partial\nabla$.

\begin{gather*}
  \int_{-\infty}^{\infty} e^{-x^{2}}\,dx=\sqrt{\pi},\qquad
  \prod_{k=1}^{n}\Bigl(1+\tfrac{1}{k}\Bigr)=n+1,\\
  \lim_{n\to\infty}\Bigl(1+\tfrac1n\Bigr)^{n}=e.
\end{gather*}
\begin{multline}
  \sum_{k=0}^{n}\binom{n}{k}x^{k}y^{n-k}=(x+y)^{n}
  \quad\text{and}\quad\\
  \sum_{\substack{0\le i\le n\\ i\ \text{even}}}\binom{n}{i}=2^{\,n-1}.
\end{multline}
\end{document}
